% =========================================================
% Proof Stub: Monotone Subsequence Theorem
% Source: volume-iii/analysis/sequences/notes/sequences/notes-subsequences.tex
% Mode: standard
% =========================================================

\clearpage
\phantomsection
\hypertarget{proof-RA-ABB-P-monotone-subseq}{}

\subsubsection[Monotone Subsequence Theorem]{Proof --- Monotone Subsequence Theorem}

\bigskip

\noindent
\textbf{Source.}
Tao, \emph{Analysis I}; see notes \texttt{notes-subsequences.tex}.

\vspace{0.75em}

\noindent
\textbf{Goal.}
Every sequence in $\mathbb{R}$ has a monotone subsequence.

\vspace{0.75em}

\noindent
\textbf{Logical form.}
\[
  \forall (x_n) \text{ a real sequence},\;
  \exists \text{ strictly increasing } \sigma : \mathbb{N} \to \mathbb{N}
  \text{ such that } (x_{\sigma(k)}) \text{ is monotone.}
\]

\vspace{0.75em}

\noindent
\textbf{Proof strategy.}
Call $n$ a \emph{peak index} if $x_n \ge x_m$ for all $m \ge n$.
\textbf{Case 1:} There are infinitely many peak indices $n_1 < n_2 < \cdots$.
Then $(x_{n_k})$ is decreasing.
\textbf{Case 2:} There are only finitely many peak indices. Then beyond
the last peak, for each $n$ there exists $m > n$ with $x_m > x_n$; this
allows construction of a strictly increasing subsequence.

\vspace{0.75em}

\noindent
\textbf{Note.}
This is a key lemma that makes Bolzano-Weierstrass short.
The peak-index construction is elegant and worth understanding deeply.

\vspace{0.75em}

\noindent
\textbf{Proof.}
\begin{proof}
\Claim Every real sequence has a monotone subsequence.

\Given $(x_n)$ a real sequence.

\Goal Construct a monotone subsequence.

\Strategy Define peak indices; either infinitely many (giving decreasing
subsequence) or finitely many (giving increasing subsequence).

\medskip

% --- define: n is a peak index if ∀ m ≥ n, x_n ≥ x_m ---

\textbf{Case 1: Infinitely many peak indices.}
% --- let n_1 < n_2 < ··· be the peak indices ---
% --- show x_{n_1} ≥ x_{n_2} ≥ ··· (decreasing) ---

\medskip
\textbf{Case 2: Finitely many peak indices.}
% --- let N be beyond all peak indices ---
% --- construct increasing subsequence: start at any n_1 ≥ N ---
% --- since n_1 is not a peak: ∃ n_2 > n_1 with x_{n_2} > x_{n_1} ---
% --- iterate to get strictly increasing subsequence ---

\end{proof}

\begin{remark*}[Proof shape]
Case split on finiteness of peak indices. In each case the subsequence
is constructed explicitly.
\end{remark*}

\begin{remark*}[Dependencies]
\begin{itemize}
  \item Well-ordering of $\mathbb{N}$ (to enumerate peak indices).
  \item Definition of monotone and subsequence.
\end{itemize}
\end{remark*}
